\begin{explanationbox}
	\textbf{\textcolor{CExplanation}{一句话直觉}}

	轴对称的本质是镜面反射：两点关于直线对称，相当于平面沿该直线折叠后两点重合。因此对称轴就是两点连线的垂直平分线。

	\medskip
	\textbf{\textcolor{CExplanation}{核心拆解}}
	\begin{description}[style=nextline, leftmargin=2.8em, labelsep=0.5em, itemsep=0.45em]
		\item[\textbf{要点一（垂直关系）}] 对称轴与对称点连线互相垂直。
		\item[\textbf{要点二（中点关系）}] 对称轴经过对称点连线的中点。
	\end{description}

	\medskip
	\textbf{\textcolor{CExplanation}{几何本质（对称即反射）}}

	沿直线 $l$ 反射（折叠）时，$l$ 上的点不动，$l$ 两侧的图形互为镜像。因此 $l$ 必然垂直平分连接镜像点的线段。

	\medskip
	\textbf{\textcolor{CExplanation}{代数意义（坐标表达）}}

	在解析几何中，若 $P(x_1,y_1)$ 与 $P'(x_2,y_2)$ 关于直线 $l: ax+by+c=0$ 对称，则中点 $\left(\frac{x_1+x_2}{2},\frac{y_1+y_2}{2}\right)$ 在 $l$ 上，且 $k_{PP'}\cdot k_l = -1$（若斜率存在）。

	\medskip
	\textbf{\textcolor{CExplanation}{考点价值（高频应用）}}
	\begin{itemize}[leftmargin=*, itemsep=0.5em, topsep=0.4em]
		\item \textbf{考法一（求对称点或轴）} 利用中点和垂直列方程组求解。
		\item \textbf{考法二（证明对称关系）} 验证中点与垂直，或利用对称性转移线段、角度。
		\item \textbf{考法三（构造对称求最值）} 将军饮马型问题，通过对称化折线为直线。
	\end{itemize}

	\medskip
	\textbf{\textcolor{CExplanation}{顿悟点}}

	\textbf{对称轴就是两点之间的“镜子”：它既平分又垂直，两个条件缺一不可。}

	\medskip
	\textbf{\textcolor{CExplanation}{使用场景}}
	\begin{itemize}[leftmargin=*, itemsep=0.5em, topsep=0.4em]
		\item \textbf{场景一（翻折问题）} 平面几何折叠后重合，直接推出中垂线。
		\item \textbf{场景二（对称建模）} 在函数与解析几何中，利用对称简化距离或角度运算。
	\end{itemize}
\end{explanationbox}
